\documentclass[10pt,a4paper]{amsart}
\usepackage[margin=19mm]{geometry}
\usepackage{amsmath,amssymb,mathtools}
\usepackage[scr=rsfs,cal=euler]{mathalfa}
\usepackage{xfrac}
\usepackage[unicode,hidelinks,bookmarks=false]{hyperref}
\newcommand{\ie}{i.e.\ }
\newcommand{\cf}{cf.\ }
\newcommand{\resp}{respectively\ }
\newcommand{\Subsection}{\subsection}
\providecommand{\nobreakdash}{\nobreak\hbox{-}}
\setlength{\emergencystretch}{2em}
\multlinegap0pt
\usepackage{amssymb}
\usepackage{mathtools}
\usepackage{float}
\usepackage{xfrac}
\newtheorem{defn}{Definition}
\newtheorem{prop}{Proposition}
\newtheorem{lem}{Lemma}
\newtheorem{thm}{Theorem}
\newcommand\ed{\operatorname{End}}
\newcommand{\f}{\varphi}
\newcommand{\g}{\psi}
\title{On the binary relations defined using GD1 and 1GD inverses over infinite dimensional vector spaces}
\author{Diego Alba Alonso, Javier S\'anchez Gonz\'alez}
\date{Source: arXiv:2604.26723v1 (CC BY 4.0)}
\begin{document}
\maketitle

\noindent\emph{Source and licence.} On the binary relations defined using GD1 and 1GD inverses over infinite dimensional vector spaces. Version arXiv:2604.26723v1 (CC BY 4.0), distributed by arXiv under the Creative Commons Attribution 4.0 International licence, http://creativecommons.org/licenses/by/4.0/, as recorded in the arXiv licence metadata for this exact version and shown on the abstract page https://arxiv.org/abs/2604.26723v1. Full licence notice: \texttt{licenses/arxiv-2604.26723-CC-BY-4.0.txt}.\par\smallskip

\noindent\textit{Editorial scope.} Counted unit: the article's thm equiv (source0.tex lines 980--987). The article's complete original proof of it is delivered with it (source0.tex lines 988--1003, one \texttt{proof} environment of 16 lines). No mathematics is written, repaired or re-derived: the statement and the proof are the article's own lines, in the article's own notation, and no retained line is substituted or re-flowed.  Retained uncounted as the source-local context the proof needs: lines 349--356; lines 385--387; lines 393--395.  Nothing was omitted from any retained span, so the package carries no omission record.  No retained line contains a citation, so the wrapper carries no bibliography and no bibliography entry.  No retained line contains a figure environment, an image-inclusion command or drawing code, so no image is delivered and the package declares no asset dependency.  Editorial formatting only, in addition: an AMS \texttt{amsart} preamble that restates the article's own declarations and macros used by the retained lines, namely the environments \texttt{defn}, \texttt{prop}, \texttt{lem}, \texttt{thm} on separate counters, together with the standard packages those lines need.  The retained environments print in this document as Definition 1, Proposition 1, Lemma 1, Theorem 1 in order of appearance, whereas the article prints the same items with the numbering of its own full document.  The complete native source of the article as distributed in the arXiv e-print for this exact version is bundled unchanged as \texttt{sources/199414/source0.tex}.
\begin{defn}\label{D: GD1 }
Let $\f \in \ed_k(V)$ be a finite potent endomorphism. We will call ``generalized Drazin 1 inverse'' of $\f,$ or simply by $GD1$ inverse, to any linear operator, denoted as $\f^{GD1}$, satisfying that: \begin{align*}
\f^{GD1}\circ \f \circ \f^{GD1}& =\f^{GD1}\\
\f^{GD1}\circ \f & = \f^{GD}\circ \f \\
\f\circ \f^{GD1} & = \f \circ \f^-,
\end{align*} 
for certain $\f^-\in X_{\f}(1)$ and $\f^{GD}\in X_{\f}(GD).$
\end{defn} 

\begin{prop}\label{T: Charact Algebraica}
Let $\f \in \ed_k(V)$ be a finite potent endomorphism. Then, an endomorphism $\g\in \ed_k(V)$ satisfies that $\g=\f^{GD1}$ (in the sense of Definition \ref{D: GD1 }) if and only if $\g=\f^{GD}\circ \f\circ \f^{-}$ for certain $\f^-\in X_{\f}(1)$ and $\f^{GD}\in X_{\f}(GD).$ 
\end{prop}

\begin{lem}\label{L: GD1 son reflex}
Let $\f \in \ed_k(V)$ be a finite potent endomorphism. Then, every GD1 inverse is a reflexive generalized inverse.
\end{lem}

\begin{thm}\label{T:3 equiv  }
Let $\f,\g \in \ed_k(V)$ be two finite potent endomorphisms. Then, the following statements are equivalent:
\begin{itemize}
\item[I.)] $\f \, \leq^{GD1} \, \g$  ;
\item[II.)] $\f=\f\circ \f^{-}\circ \g=\g \circ \f^{GD}\circ \f$ for certain $\f^{-}\in X_{\f}(1)$ and $\f^{GD}\in X_{\f}(GD);$  
\item[III.)] $\f=\f\circ \f^{GD1}\circ \g=\g\circ \f^{GD1}\circ \f$ for certain $\f^{GD1}\in X_{\f}(GD1).$
\end{itemize}
\end{thm}

\begin{proof}
$I) \Longrightarrow II).$ As $\f \circ \f^{GD1}=\g\circ \f^{GD1}$ and $\f^{GD1}\circ \f =\f^{GD1}\circ \g $ for certain $\f^{GD1}\in X_{\f}(GD1).$ In virtue of the expression of any GD1 inverse obtained in Proposition \ref{T: Charact Algebraica}, we know that $\f\circ \f^{GD}\circ \f \circ \f^{-}=\g\circ\f^{GD}\circ \f \circ \f^{-} $ and therefore \begin{equation}\label{eq: 3equiv a}
\f\circ \f^{-}=\g\circ \f^{GD}\circ \f \circ \f^{-}.
\end{equation}
Similarly, $\f^{GD}\circ \f \circ \f^{-}\circ \f =\f^{GD}\circ \f \circ \f^{-}\circ \g ,$ so \begin{equation}\label{eq: 3equiv b}
\f^{GD}\circ \f=\f^{GD}\circ \f \circ \f^{-}\circ \g .
\end{equation}
Composing in \eqref{eq: 3equiv a} with $\f$ on the right hand side and in \eqref{eq: 3equiv b} with $\f$ on the left hand side, one obtains that: $\f=\g\circ \f^{GD}\circ \f$ and $\f=\f\circ\f^{-}\circ \g$ respectively, so the first implication is proved. \\
$II) \Longrightarrow III).$ Let us suppose that $\f=\f\circ \f^{-}\circ \g=\g \circ \f^{GD}\circ \f$ holds for certain $\f^{GD}\in X_{\f}(GD)$ and $\f^{-}\in X_{\f}(1).$ Take the composition: $\f^{GD}\circ \f \circ \f^{-}$ which is a GD1 inverse of $\f .$ By direct computation, one gets that: \begin{align*}
\f\circ (\f^{GD}\circ \f \circ \f^{-})\circ \g & = \f \circ \f^{-}\circ \g =\f ;\\
\g\circ (\f^{GD}\circ \f \circ \f^{-})\circ \f & =\g \circ \f^{GD}\circ \f=\f ,
\end{align*}
so the condition in the statement is deduced.\\
$III)\Longrightarrow I).$ Let us now suppose that $\f=\f\circ \f^{GD1}\circ \g=\g\circ \f^{GD1}\circ \f $ for certain $\f^{GD1}.$ To wit, right-composing with $\f^{GD1}$ in the first equality one gets: $$\f\circ \f^{GD1}=\g \circ \f^{GD1}\circ \f \circ \f^{GD1}=\g \circ \f^{GD1} $$ 
as GD1 inverses are, in particular, $\{2 \}-$inverses (recall Lemma \ref{L: GD1 son reflex}). Reasoning on an analogous way, one obtains that: $$\f^{GD1}\circ \f=\f^{GD1}\circ \f \circ \f^{GD1}\circ \g =\f^{GD1}\circ \g $$ so both conditions of the relation $\leq^{GD1}$ hold and hence the claim is proved.
\end{proof}

\end{document}
